\documentclass[11pt]{article}

\usepackage[a4paper,margin=2.4cm]{geometry}
\usepackage{lmodern}
\usepackage[T1]{fontenc}
\usepackage[utf8]{inputenc}
\usepackage{microtype}
\usepackage{graphicx}
\usepackage{xcolor}
\usepackage{float}
\usepackage{caption}
\usepackage{enumitem}
\usepackage{parskip}
\usepackage{titlesec}
\usepackage[hidelinks]{hyperref}

% Screenshots are the same files the in-app help page uses.
\graphicspath{{../../frontend/public/help/}}

\definecolor{accent}{HTML}{1D4ED8}
\definecolor{accentbg}{HTML}{EFF4FF}
\definecolor{soft}{HTML}{555555}

\titleformat{\section}{\Large\bfseries\color{accent}}{}{0pt}{}
\titlespacing*{\section}{0pt}{1.4em}{0.5em}

\captionsetup{font={small,it},labelformat=empty,skip=4pt}
\setlength{\fboxsep}{10pt}

\newcommand{\shot}[2]{%
  \begin{figure}[H]\centering
  \includegraphics[width=\linewidth,height=0.6\textheight,keepaspectratio]{#1}%
  \caption*{#2}%
  \end{figure}}

\newcommand{\callout}[1]{%
  \begin{center}
  \fcolorbox{accent}{accentbg}{\parbox{0.96\linewidth}{\small #1}}
  \end{center}}

\hypersetup{
  pdftitle={Minerva: Teacher Setup Guide},
  pdfauthor={DSV, Stockholm University},
}

\begin{document}

\begin{center}
  {\LARGE\bfseries Minerva}\\[2pt]
  {\Large\color{accent} Teacher Setup Guide}\\[6pt]
  {\color{soft}\small Connect the Minerva AI study assistant to your Moodle course \textbullet{} about 5 minutes \textbullet{} 3 steps}
\end{center}

\vspace{0.4em}
\hrule height 1pt
\vspace{1.0em}

\noindent
Minerva is an AI study assistant that answers from your own course materials.
Your course, its materials and its recorded lectures are already kept up to
date in Minerva, so all that is left is putting a link to it inside Moodle.

\callout{\textbf{Before you start:} your administrator has registered Minerva as
an external tool across Moodle. You still need to enable it for each of your
courses (Step 1) before it shows up in the activity chooser. If you do not see
the options shown here, contact \href{mailto:lambda@dsv.su.se}{lambda@dsv.su.se}.}

\section{What happens on its own}

\begin{itemize}[leftmargin=1.4em,itemsep=4pt]
  \item Courses are imported from Daisy every night, so yours is already in
  Minerva.
  \item A Moodle course whose \textbf{Course ID number} holds its Daisy offering
  id is linked to the matching Minerva course automatically, usually within the
  hour.
  \item Linked material is re-synced twice an hour, so new files and pages are
  picked up without you doing anything.
  \item Recorded lectures from DSV Play are pulled in from the course's Play
  designation, transcripts and all.
\end{itemize}

\shot{minerva-dashboard.png}{Your Minerva dashboard, with the courses you teach.}

\section{Step 1. Enable Minerva in your course}

Minerva is registered across Moodle, but it does not appear in a course until
you switch it on. In your Moodle course, open the \textbf{More} menu and choose
\textbf{LTI External tools}. Find \textbf{Minerva}, activate it, and turn on
\textbf{Show in activity chooser}.

\shot{enable-tool.png}{Course LTI External tools: activate Minerva and show it in the activity chooser.}

Once enabled correctly, Minerva appears in the activity chooser with its logo.
If you instead see a blank or generic icon, the tool was registered without its
icon; contact \href{mailto:lambda@dsv.su.se}{lambda@dsv.su.se}.

\section{Step 2. Add the Minerva activity to a section}

In your Moodle course, turn on \textbf{Edit mode}, select \textbf{Add an
activity or resource} in the section where you want the assistant, and choose
\textbf{Minerva}.

\shot{activity-chooser.png}{Pick Minerva from the activity chooser; it carries the Minerva logo.}

Give it a name students will recognise (for example, ``Minerva AI assistant''),
then choose \textbf{Save and return to course}. The activity now appears in your
course.

\shot{lti-config.png}{Name the activity and save.}
\shot{course-activity.png}{The Minerva activity is now in your course.}

\section{Step 3. Open it once and pick the course}

Click the Minerva activity you added. The first time, pick which Minerva course
it should open and choose \textbf{Link and open Minerva}. Minerva remembers
this, so every later launch by you or your students opens the right course.

\shot{first-launch-bind.png}{On first launch, connect the activity to your Minerva course.}

\section{That is it: here is what students see}

Students click the same activity and get a chat assistant that answers from your
course materials, right inside Moodle. The suggested questions are drawn from
those materials.

\shot{student-chat.png}{The student view: a course-grounded chat assistant.}
\shot{launch-in-moodle.png}{The same assistant, embedded inside the Moodle activity.}

\section{Deciding exactly what is indexed}

The automatic link indexes what is in your Moodle course. If you would rather
choose precisely what the assistant answers from, create your own course in
Minerva with \textbf{New Course} and upload only the documents you want.

Then add a second Minerva activity in Moodle and point it at that course the
first time you open it. Both activities can sit in the same Moodle course;
students get whichever one you send them to.

\section{If the automatic link did not happen}

Without a Daisy offering id in the Moodle course's \textbf{Course ID number}
there is nothing to match on, so nothing gets linked. Open \textbf{Minerva
settings} from the course's \textbf{More} menu, pick the Minerva course and
choose \textbf{Link Minerva course}.

\shot{link-form.png}{Minerva settings: choose your Minerva course and link it.}

\textbf{Sync materials} then uploads your files, pages and section summaries
right away; after that the twice-hourly sync keeps them current.

\shot{linked.png}{Once linked, you can sync materials or unlink at any time.}
\shot{sync-complete.png}{Minerva confirms how many resources were uploaded.}

\section{Good to know}

\begin{itemize}[leftmargin=1.4em,itemsep=4pt]
  \item The assistant only answers from material in the Minerva course the
  activity points at. Student conversations are visible to course staff; see
  \textbf{Data handling} in Minerva for details.
  \item A course linked automatically says so in \textbf{Minerva settings}:
  unlinking it there keeps it unlinked, and it will not be auto-linked again.
  \item Teacher-answered forum threads can be included from \textbf{Minerva
  settings} (optional; student names are removed before upload).
  \item Questions or need higher usage limits? Email
  \href{mailto:lambda@dsv.su.se}{lambda@dsv.su.se}.
\end{itemize}

\vspace{1.2em}
\hrule height 0.4pt
\vspace{0.4em}
\noindent{\color{soft}\small Minerva \textbullet{} Department of Computer and
Systems Sciences (DSV), Stockholm University \textbullet{}
\href{mailto:lambda@dsv.su.se}{lambda@dsv.su.se}}

\end{document}
